% sections/model.tex — The Persistent GPU Actor Model
\section{The Persistent GPU Actor Model}
\label{sec:model}

This section develops the formal model.  We define the syntactic
domain of actors and configurations~(\Cref{sec:model:syntax}), then
give a small-step operational
semantics for the lifecycle~(\Cref{sec:model:lifecycle}), and
finally state and prove the basic safety property:
no actor reaches \Active{} without first having been
initialized~(\Cref{thm:lifecycle-soundness}).
The mechanized \TLA{} specification corresponding to
\Cref{sec:model:lifecycle} ships with the open-source distribution
of \system{} as
\texttt{docs/verification/actor\_lifecycle.tla} and has been
model-checked under the configuration
\texttt{actor\_lifecycle.cfg} for finite parameter ranges.

\subsection{Syntactic Domain}
\label{sec:model:syntax}

We work over a countable set of \emph{actor identifiers}
\(\mathcal{A} = \{a_1, a_2, \ldots\}\) with a distinguished
root \(\rho \in \mathcal{A}\).  Each actor has an associated
\emph{parent} given by a function
\(\pi : \mathcal{A} \to \mathcal{A} \cup \{\rho\}\); the parent of
the root is itself.  We assume the parent relation forms a tree
(no cycles; every non-root actor has a finite ancestor chain
ending at \(\rho\)).

\begin{definition}[Lifecycle State]
\label{def:lifecycle-state}
The set of \emph{lifecycle states} is
\[
\Sigma_{\mathrm{life}} \;=\;
  \{\Dormant{}, \Initializing{}, \Active{},
    \Draining{}, \Terminated{}, \Failed{}\}.
\]
\Dormant{} is the slot-empty state of an unallocated supervisor
entry; \Failed{} is a sink for unrecovered faults that the
supervisor has not yet acted upon.
\end{definition}

\begin{definition}[Actor State Vector]
\label{def:actor-state}
An actor's per-actor state is a tuple
\[
  \sigma \;=\;
  \langle \mathit{life},\,
          \mathit{snap},\,
          \mathit{live},\,
          \mathit{hlc},\,
          \mathit{rcount},\,
          \mathit{tag} \rangle
\]
where
\(\mathit{life} \in \Sigma_{\mathrm{life}}\) is the lifecycle state;
\(\mathit{snap} \in \mathit{Bytes}^*\) is the bytes of the most
recent state snapshot (empty when no snapshot has been taken);
\(\mathit{live} \in \mathit{Bytes}^*\) is the bytes of the actor's
live in-memory state, evolving while \(\mathit{life} = \Active{}\);
\(\mathit{hlc} \in \N^2\) is the actor's Hybrid Logical Clock
(physical-counter pair);
\(\mathit{rcount} \in \N\) is the supervisor's restart counter for
this actor within the current restart window; and
\(\mathit{tag} \in \N \times \N\) is the audit boundary
\((\mathit{org\_id}, \mathit{engagement\_id})\) under which the
actor is registered.
\end{definition}

\begin{definition}[Configuration]
\label{def:configuration}
A \emph{configuration} \(C\) of the actor system is a triple
\(C = \langle \Theta, M, t \rangle\), where
\(\Theta : \mathcal{A} \to (\mathit{Defined} \cup \{\bot\})\) is
the partial state assignment
(\(\Theta(a) = \bot\) means the actor identifier \(a\) is
unallocated), \(M\) is the K2K message store
(\Cref{sec:k2k}), and
\(t \in \N\) is the global step counter.
We write \(\mathcal{C}\) for the set of all configurations.
\end{definition}

The state assignment \(\Theta\) is finite-support: only finitely
many identifiers map to defined states at any one time.

\subsection{Lifecycle Transitions}
\label{sec:model:lifecycle}

The lifecycle of an actor is governed by six rules: spawn, activate,
process, quiesce, terminate, and fail.  Each is given as a
small-step transition over configurations.  We use the notation
\(C \trans{r} C'\) to mean ``configuration \(C\) takes one step under
rule~\(r\) to configuration \(C'\).''  \Cref{fig:lifecycle} summarizes
the resulting state machine.

\begin{figure}[h]
\centering
\begin{tikzpicture}[
    node distance=15mm and 22mm,
    every state/.style={draw, rounded corners=2pt, minimum size=10mm,
                        font=\small\sffamily, inner sep=2pt},
    every edge/.style={draw, ->, >=Stealth, thin},
    every node/.style={font=\footnotesize\sffamily}
  ]
  \node[state] (D)  {Dormant};
  \node[state, right=of D]  (I)  {Initializing};
  \node[state, right=of I]  (A)  {Active};
  \node[state, right=of A]  (Dr) {Draining};
  \node[state, right=of Dr] (T)  {Terminated};
  \node[state, below=12mm of A, fill=red!8] (F) {Failed};

  \path (D)  edge node[above]{\textsc{Spawn}}     (I);
  \path (I)  edge node[above]{\textsc{Activate}}  (A);
  \path (A)  edge node[above]{\textsc{Quiesce}}   (Dr);
  \path (Dr) edge node[above]{\textsc{Terminate}} (T);
  \path (T)  edge[bend left=35] node[below]{\textsc{Release}} (D);

  \path (A)  edge[loop above] node{\textsc{Process}} (A);
  \path (I)  edge[bend left=20] node[left, near end]{\textsc{Fail}} (F);
  \path (A)  edge node[right]{\textsc{Fail}} (F);
  \path (Dr) edge[bend right=20] node[right, near end]{\textsc{Fail}} (F);
  \path (F)  edge[bend left=40] node[below left]{\textsc{Restart}} (I);
\end{tikzpicture}
\caption{Lifecycle state machine.  Solid edges are SOS rules
(\Cref{rule:spawn,rule:activate,rule:process} and the appendix
rules).  \Failed{} is a fault attractor reachable from any
non-terminal state; \textsc{Restart}~(\Cref{rule:restart}) returns
to \Initializing{} via the snapshot-restore protocol of
\Cref{sec:supervision:restart}.}
\label{fig:lifecycle}
\end{figure}

\paragraph{Spawn}
A previously-unallocated identifier becomes allocated in
\Initializing{}.  The supervisor records the parent and resets the
restart counter and HLC.

\begin{equation}
\inference[\textsc{Spawn}]
{
  \Theta(a) = \bot \\
  \pi(a) = p \in \mathcal{A} \cup \{\rho\} \\
  \mathit{tag}_{\mathrm{new}} \in \N \times \N
}
{
  \langle \Theta, M, t\rangle
  \trans{\mathsf{spawn}(a)}
  \langle \Theta', M, t+1\rangle
}
\label{rule:spawn}
\end{equation}

\noindent where
\(\Theta' = \Theta[a \mapsto
  \langle \Initializing{}, \varepsilon, \varepsilon, \langle 0, 0\rangle,
          0, \mathit{tag}_{\mathrm{new}}\rangle]\)
and \(\varepsilon\) is the empty byte sequence.

\paragraph{Activate}
An actor in \Initializing{} transitions to \Active{}.  The
supervisor takes a snapshot of the current live state (which is
\(\varepsilon\) on the first activation but populated on restart).

\begin{equation}
\inference[\textsc{Activate}]
{
  \Theta(a) = \langle \Initializing{}, \mathit{snap}, \mathit{live},
                       h, r, \tau\rangle
}
{
  \langle \Theta, M, t\rangle
  \trans{\mathsf{activate}(a)}
  \langle \Theta', M, t+1\rangle
}
\label{rule:activate}
\end{equation}

\noindent where
\(\Theta' = \Theta[a \mapsto
  \langle \Active{}, \mathit{live}, \mathit{live}, h, r, \tau\rangle]\).

The snapshot is set to \(\mathit{live}\) at activation: this is the
\emph{snapshot invariant}—\(\mathit{snap}\) is always equal to the
\(\mathit{live}\) bytes at the most recent activation point.

\paragraph{Process}
An \Active{} actor advances its live state.  The transition
function \(\delta : \mathit{Bytes}^* \times M \to
\mathit{Bytes}^* \times M\) consumes a finite prefix of the actor's
inbox in \(M\), updates \(\mathit{live}\), and may produce new
messages.

\begin{equation}
\inference[\textsc{Process}]
{
  \Theta(a) = \langle \Active{}, \mathit{snap}, \mathit{live},
                       h, r, \tau\rangle \\
  \delta(\mathit{live}, M) = (\mathit{live}', M')
}
{
  \langle \Theta, M, t\rangle
  \trans{\mathsf{process}(a)}
  \langle \Theta', M', t+1\rangle
}
\label{rule:process}
\end{equation}

\noindent where
\(\Theta' = \Theta[a \mapsto
  \langle \Active{}, \mathit{snap}, \mathit{live}', \hlc.\mathsf{tick}(h), r, \tau\rangle]\)
and \(\hlc.\mathsf{tick}\) is the HLC advance operation defined in
\Cref{sec:supervision}.

\paragraph{Quiesce, Terminate, Fail}
\Cref{rule:quiesce,rule:terminate,rule:fail} (omitted here for space;
see \Cref{app:lifecycle-rules}) give the remaining transitions.
\Quiesce{} moves an actor from \Active{} into \Draining{} so its
inbox can flush; \Terminate{} releases the slot, returning to
\Dormant{}; \Fail{} models a fault that takes the actor to
\Failed{}, where the supervisor decides whether to restart it
(per \Cref{sec:supervision}) or destroy it.

\subsection{Soundness}
\label{sec:model:soundness}

We can now state the basic safety property: no actor takes a
\textsc{Process} step from a state in which it has not been
activated.  In informal terms: a kernel never executes user code
before its initial state has been laid down.

\begin{theorem}[Lifecycle Soundness]
\label{thm:lifecycle-soundness}
For all configurations \(C\) reachable from any initial
configuration~\(C_0\) under the rules of
\Cref{sec:model:lifecycle}, and for every actor identifier \(a\)
with \(\Theta(a) \ne \bot\) in \(C\):
\[
  \text{\textsc{Process}}(a) \text{ has been applied along the trace}
  \;\Longrightarrow\;
  \text{\textsc{Activate}}(a) \text{ was applied earlier in the trace.}
\]
\end{theorem}

\begin{proof}[Proof sketch]
By induction on the length of the reduction sequence
\(C_0 \trans{r_1} C_1 \trans{r_2} \cdots \trans{r_n} C\).
The base case \(n = 0\) is vacuous because no \textsc{Process}
step has been taken.  For the induction step, examine the rule
\(r_n\) applied at step \(n\).  By inspection of \Cref{rule:process},
\textsc{Process}(\(a\)) requires
\(\Theta_{n-1}(a) = \langle \Active{}, \ldots\rangle\).  By
\Cref{def:lifecycle-state}, the only rule whose target state is
\Active{} is \textsc{Activate}~(\Cref{rule:activate}).  By the
induction hypothesis applied to the prefix
\(C_0 \trans{r_1} \cdots \trans{r_{n-1}} C_{n-1}\), an
\textsc{Activate}(\(a\)) step exists in that prefix; this completes
the induction.
\qed
\end{proof}

\Cref{thm:lifecycle-soundness} is the analogue of the classical
``no use before initialization'' property of imperative type
systems~\citep{wright1994syntactic}.  It is also a precondition for
the more substantive theorems of \Cref{sec:supervision}, which
establish that snapshot-restore preserves causal time even across a
fault.

\subsection{Mechanization}
\label{sec:model:mechanization}

\Cref{sec:model:lifecycle} is mechanized in \TLA{} in
\texttt{docs/verification/actor\_lifecycle.tla}.  The mechanization
is faithful to the operational semantics: every rule corresponds to
a \TLA{} action, and the variables \(\mathit{state},
\mathit{snapshot}, \mathit{live\_state}\) match the components of
\Cref{def:actor-state}.  We model-checked the specification under
\texttt{actor\_lifecycle.cfg} for parameter ranges
\(|\mathcal{A}| \le 5\) and \(\mathit{MaxSteps} \le 30\); within
those bounds, no violation of \Cref{thm:lifecycle-soundness} or its
companion invariants (snapshot-equals-live-at-activation; restart
preserves snapshot) was observed across the explored state space of
\(\sim 10^7\) distinct configurations.  The mechanized proof is the
empirical floor of the analytical proof above.
